\section{The BNT canonical form on a sector decomposition}
\label{sec:sector_bnt_api}

% Source: Cirac2016MPDO_arXiv \S~II lines 287--301; Cirac2021Matrix
% Definition~4.2, Proposition~4.3, lines 1864--1884.
This section introduces the BNT canonical-form conditions of
\cite[Section~2]{Cirac2016MPDO_arXiv}, defined on a sector decomposition with raw
sector weights $\mu_{j,q}$ and coefficients $c_N^{(j)} = \sum_q \mu_{j,q}^{\,N}$;
no equal-modulus factorization is assumed.

\subsection{Sector decompositions and phase matching}
\label{sec:bnt_sector_decompositions}

The sector decomposition and the MPV phase-equivalence relations introduced
here determine the BNT canonical-form conditions below.

\begin{definition}[Sector weights and sector decomposition]%
    \label{def:sector_weight_data}
    \lean{MPSTensor.SectorWeightData, MPSTensor.SectorDecomposition}
    \lean{MPSTensor.SectorDecomposition.toTensor}
    \lean{MPSTensor.SectorDecomposition.totalDim}
    \lean{MPSTensor.SectorDecomposition.flatDim_flatIndexEquiv,
        MPSTensor.SectorDecomposition.flatWeight_flatIndexEquiv,
        MPSTensor.SectorDecomposition.flatBasis_flatIndexEquiv_heq}
    \leanok
    A sector-weight tuple over $g$ basis blocks consists of multiplicities
    $r_j > 0$, nonzero weights $\mu_{j,q}$ for
    $q \in \{1,\ldots,r_j\}$, and the sector coefficients
    \begin{align}
        c_{S,j}^{(N)}
         & :=
        \sum_{q=1}^{r_j} \mu_{j,q}^N.
        \label{eq:bnt_sector_coeff_def}
    \end{align}
    A \emph{sector decomposition} consists of the basis blocks
    $(A_j)_{j=1}^g$ together with these multiplicities and weights. If $A_j$
    has bond dimension $D_j$, then the flattened copy family has total bond
    dimension
    \begin{align}
        D_{\mathrm{tot}}
         & =\sum_{j=1}^{g}r_jD_j.
        \notag
    \end{align}
\end{definition}

\begin{definition}[Physical blocking of a sector decomposition]%
    \label{def:sector_decomposition_block_tensor}
    \lean{MPSTensor.SectorDecomposition.blockTensor}
    \leanok
    \uses{def:sector_weight_data, def:blocked_tensor}
    Let $P$ be a sector decomposition with representatives $A_j$,
    multiplicities $r_j$, and copy weights $\mu_{j,q}$. Its physical
    $p$-site blocking has the same representative and copy indices, with
    \begin{align}
        A_j       & \longmapsto A_j^{[p]},
        \notag                                 \\
        \mu_{j,q} & \longmapsto \mu_{j,q}^{p}.
        \notag
    \end{align}
    The two-layer decomposition is
    \cite[lines~283--301]{Cirac2016MPDO_arXiv}. Blocking preserves the copy
    index sets and bond dimensions, replaces each flattened representative by
    its $p$-site blocking, and raises each flattened weight to the $p$-th power.
\end{definition}

\begin{theorem}[Indices, dimensions, weights, and representatives under blocking]%
    \label{thm:sector_decomposition_block_tensor_data}
    \lean{MPSTensor.SectorDecomposition.blockTensor_basis}
    \lean{MPSTensor.SectorDecomposition.blockTensor_weight}
    \lean{MPSTensor.SectorDecomposition.blockTensor_copies}
    \lean{MPSTensor.SectorDecomposition.blockTensor_totalCopies}
    \lean{MPSTensor.SectorDecomposition.blockTensor_flatDim}
    \lean{MPSTensor.SectorDecomposition.blockTensor_totalDim}
    \lean{MPSTensor.SectorDecomposition.blockTensor_flatWeight}
    \lean{MPSTensor.SectorDecomposition.blockTensor_flatBasis}
    \leanok
    \uses{def:sector_decomposition_block_tensor}
    Physical blocking preserves the copy sets and bond dimensions. On the
    flattened copy index, it raises the weight to the $p$-th power and replaces
    the representative by its $p$-site blocking.
\end{theorem}

\begin{proof}\leanok
    \uses{def:sector_decomposition_block_tensor}
    These identities follow directly from the definition.
\end{proof}

\begin{lemma}[Sector decomposition formula]%
    \label{lem:paper_decomp_formula}
    \lean{MPSTensor.SectorDecomposition.mpv_toTensor_eq_sum_coeff}
    \leanok
    \uses{def:sector_weight_data}
    For a sector decomposition $S$ with associated tensor $A_{\mathrm{tot}}$,
    the MPV satisfies
    \begin{align}
        V^{(N)}(A_{\mathrm{tot}})_\sigma
         & =
        \sum_{j=1}^{g} c_{S,j}^{(N)} V^{(N)}(A_j)_\sigma.
        \label{eq:bnt_sector_decomp}
    \end{align}
\end{lemma}

Two blocks are identified for canonical-form purposes when their MPV families
differ by a scalar power. The next three definitions formalize this relation,
the common cover it generates between two families, and the phase classes of a
single family.

\begin{definition}[MPV phase equivalence for nonzero-weight blocks]
    \label{def:mpv_block_phase_equiv}
    \lean{MPSTensor.MPVBlockPhaseEquiv}
    \leanok
    \uses{def:mpv}
    Two blocks $A$ and $B$ are MPV-phase equivalent if there is a nonzero
    scalar $\zeta$ such that, for every positive system size~$N$ and every
    physical word, $V^{(N)}(B)_\sigma=\zeta^N V^{(N)}(A)_\sigma$.
    The two bond dimensions are allowed to differ.
\end{definition}

\begin{definition}[Common MPV phase cover]
    \label{def:mpv_common_phase_cover}
    \lean{MPSTensor.MPVCommonPhaseCover}
    \leanok
    \uses{def:mpv_block_phase_equiv}
    For two finite nonzero-weight block families, a common MPV phase cover consists of a third
    finite family of blocks, a map from each nonzero-weight block family onto that common
    family, and MPV-phase equivalences between each block and its common
    representative. Both maps are required to be surjective, so the common family
    contains no unused representative.
\end{definition}

\begin{definition}[MPV phase classes]%
    \label{def:mpv_phase_class_data}
    \lean{MPSTensor.MPVPhaseEquiv}
    \lean{MPSTensor.MPVPhaseClassData}
    \lean{MPSTensor.mpvPhaseClassData}
    \lean{MPSTensor.mpvPhaseClassData_enum_zero_eq_repr}
    \lean{MPSTensor.MPVPhaseClassData.enumEquiv}
    \leanok
    For a finite family of blocks, put $j \sim k$ when there is a nonzero
    scalar factor $\zeta$ such that
    $\ket{V^{(N)}(B_k)}=\zeta^N\ket{V^{(N)}(B_j)}$ for every positive
    length~$N$. The associated tuple enumerates the finite quotient,
    chooses one representative per class, enumerates that representative first
    in its class, records the scalar relating the representative to each member,
    and proves that distinct representatives are not MPV-phase equivalent. The
    pairs $(j,q)$, with $j$ a phase class and $q$ a copy in that class, are in
    bijection with the original block indices.
\end{definition}

\begin{theorem}[Active-block BNT existence for CPSV canonical form]
    \label{thm:cpsv_canonical_form_bnt_existence}
    \lean{MPSTensor.CPSVCanonicalFormData.exists_isCPSVBasisOfNormalTensors}
    \leanok
    \uses{def:cpsv_canonical_form, def:bnt_cpsv}
    % Maintainer source anchors: CPSV16 prop:char-BNT, lines 271--280 and 1137--1146.
    Every tensor with literal CPSV canonical-form data admits a finite basis of
    normal tensors $(B_j)_{j=1}^{g}$, each of positive bond dimension. The
    family is obtained from
    $K_{\mathrm{act}}=\{k:\mu_k\ne0\}$ by choosing one representative from
    each matrix-product-vector phase class of the active blocks.

    This statement does not cover zero-weight listed blocks: they contribute
    nothing to any positive-length matrix product vector. It also makes no
    unrestricted uniqueness claim about arbitrary bases of normal tensors.
\end{theorem}

\begin{proof}\leanok
    \uses{thm:cpsv_canonical_form_bnt_characterization, def:mpv_phase_class_data, thm:cpsv_normal_mpv_phase_gauge}
    Discard the zero-weight blocks and partition the remaining finite family by
    matrix-product-vector phase equivalence. Choose one representative from
    each class. The representatives are normal and have positive bond
    dimensions. Distinct representatives are not gauge-phase equivalent. For
    every active block, the normal-tensor phase theorem identifies its bond
    dimension with that of its representative and realizes their phase
    relation by a similarity and a unit-modulus scalar.
    Theorem~\ref{thm:cpsv_canonical_form_bnt_characterization} now gives the
    basis-of-normal-tensors conditions.
\end{proof}

% Source anchors: CPSV16 lines 265--301, 1080--1117, and 1135--1146.
% Proposition 4.13, lines 1863--1921, is the downstream application.
\begin{theorem}[Structured active BNT refinement of CPSV canonical form]
    \label{thm:cpsv_canonical_form_active_bnt_refinement}
    \lean{MPSTensor.CPSVCanonicalFormData.ActiveBNTRefinement,
        MPSTensor.CPSVCanonicalFormData.exists_activeBNTRefinement,
        MPSTensor.CPSVCanonicalFormData.activeBNTRefinement}
    \lean{MPSTensor.CPSVCanonicalFormData.ActiveBNTRefinement.activeCopy}
    \lean{MPSTensor.CPSVCanonicalFormData.activeClassCopyEquiv,
        MPSTensor.CPSVCanonicalFormData.Inactive,
        MPSTensor.CPSVCanonicalFormData.GroupedIndex,
        MPSTensor.CPSVCanonicalFormData.groupedIndexEquiv,
        MPSTensor.CPSVCanonicalFormData.groupedCoordinateEquiv,
        MPSTensor.CPSVCanonicalFormData.groupedTotalDim_eq}
    \lean{MPSTensor.CPSVCanonicalFormData.ActiveBNTRefinement.groupedWeight,
        MPSTensor.CPSVCanonicalFormData.ActiveBNTRefinement.groupedBlocks,
        MPSTensor.CPSVCanonicalFormData.ActiveBNTRefinement.groupedTensor,
        MPSTensor.CPSVCanonicalFormData.ActiveBNTRefinement.regroupedTensor_eq_groupedTensor,
        MPSTensor.CPSVCanonicalFormData.ActiveBNTRefinement.groupedRegroupLetterwise,
        MPSTensor.CPSVCanonicalFormData.ActiveBNTRefinement.reconstructGrouped}
    \lean{MPSTensor.blockIndexCoordinateEquiv,
        MPSTensor.toTensorFromBlocks_eq_reindex_blockDiagonal_equiv}
    \leanok
    \uses{def:cpsv_canonical_form, def:mpv_phase_class_data, thm:cpsv_canonical_form_bnt_characterization, thm:cpsv_normal_mpv_phase_gauge, def:block_diagonal_tensor, def:global_gauge_of_blocks}
    Let
    \begin{align}
        A^i & =U^\dagger\left(\bigoplus_{k=1}^{r}\mu_kA_k^i\right)U.
        \label{eq:cpsv_active_refinement_original}
    \end{align}
    \begin{align}
        UU^\dagger & =\Id.
        \label{eq:cpsv_active_refinement_coisometry}
    \end{align}
    be literal CPSV canonical-form data. Put
    $K_{\mathrm{act}}=\{k:\mu_k\ne0\}$ and
    $K_0=\{k:\mu_k=0\}$. There are phase classes $j$, copy multiplicities
    $r_j$, and an equivalence
    \begin{align}
        \mathcal G
        :=\{(j,q):1\leq q\leq r_j\}\sqcup K_0
         & \simeq\{1,\ldots,r\}.
        \label{eq:cpsv_active_refinement_grouped_equiv}
    \end{align}
    Its active restriction identifies the class-copy pairs with
    $K_{\mathrm{act}}$. Each class has a representative $C_j$. If the copy
    $(j,q)$ corresponds to $k\in K_{\mathrm{act}}$, then its bond dimension
    equals that of $C_j$ and there are $\zeta_k\in\C$, with $|\zeta_k|=1$,
    and $X_k\in\GL_{D_k}(\C)$ such that
    \begin{align}
        A_k^i & =\zeta_kX_kC_j^iX_k^{-1}.
        \label{eq:cpsv_active_refinement_copy}
    \end{align}
    The copy retains the original coefficient $\mu_k$, and the representatives
    $(C_j)_j$ form a basis of normal tensors for $A$.

    Enumerate $\mathcal G$. The induced equivalence $\Phi$ between the
    dependent grouped block coordinates and the original flattened listed
    coordinates gives
    \begin{align}
        \sum_{x\in\mathcal G}D_x & =\sum_{k=1}^{r}D_k.
        \label{eq:cpsv_active_refinement_total_dim}
    \end{align}
    For $x=(j,q)$, let $\widehat\mu_x$ be the original weight of that copy and
    let $\widehat B_x^i=\zeta_xC_j^i$, with the required dimension
    identification. For $x\in K_0$, let $\widehat\mu_x=0$ and retain the
    original block $\widehat B_x=A_x$. Define the grouped tensor by
    \begin{align}
        T_{\mathrm{grp}}^i
         & =\operatorname{reindex}_{\Phi}
        \left(\bigoplus_{x\in\mathcal G}
        \widehat\mu_x\widehat B_x^i\right).
        \label{eq:cpsv_active_refinement_grouped_tensor}
    \end{align}
    If $G=\bigoplus_kY_k$ is the block-diagonal gauge, with $Y_k=X_k$ on
    active indices and $Y_k=\Id$ on $K_0$, then, letter by letter,
    \begin{align}
        \bigoplus_{k=1}^{r}\mu_kA_k^i
         & =GT_{\mathrm{grp}}^iG^{-1},
        \label{eq:cpsv_active_refinement_regroup}
    \end{align}
    and
    \begin{align}
        A^i & =U^\dagger GT_{\mathrm{grp}}^iG^{-1}U.
        \label{eq:cpsv_active_refinement_reconstruction}
    \end{align}
    Thus the class-copy coordinates and the inactive complement are explicit,
    while every zero-weight listed block remains present. This refinement does
    not delete the inactive coordinates; the active coisometric restriction is
    stated separately below.
\end{theorem}

\begin{proof}\leanok
    \uses{def:mpv_phase_class_data, thm:cpsv_canonical_form_bnt_characterization, thm:cpsv_normal_mpv_phase_gauge, def:block_diagonal_tensor, def:global_gauge_of_blocks}
    Enumerate the active indices by their phase classes and choose one normal
    representative in each class. The normal-tensor phase theorem gives the
    dimension equality, unit phase, and invertible gauge in
    (\ref{eq:cpsv_active_refinement_copy}). The BNT characterization proves
    that the representatives form a basis of normal tensors. Adjoin the
    inactive complement to the class-copy order. Reindexing the dependent block
    coordinates gives $\Phi$ and
    (\ref{eq:cpsv_active_refinement_total_dim}). The exact heterogeneous
    block-diagonal reindexing identity identifies the grouped tensor with the
    representative copies in the original retained-coordinate space.
    Block-diagonal conjugation then gives
    (\ref{eq:cpsv_active_refinement_regroup}); substitution into
    (\ref{eq:cpsv_active_refinement_original}) gives
    (\ref{eq:cpsv_active_refinement_reconstruction}) with the original
    coisometry $U$. In the notation of
    \cite[Proposition~4.13]{Cirac2016MPDO_arXiv}, these identities
    are
    \begin{align}
        \bigoplus_{k=1}^{r}\mu_kA_k^i
         & =GT_{\mathrm{grp}}^iG^{-1},
        \label{eq:cpsv_prop_4_13_active_bnt}            \\
        A^i
         & =U^\dagger GT_{\mathrm{grp}}^iG^{-1}U,
        \label{eq:cpsv_prop_4_13_active_reconstruction} \\
        UU^\dagger
         & =\Id.
        \label{eq:cpsv_prop_4_13_original_coisometry}
    \end{align}
\end{proof}

\begin{theorem}[Exact coisometric restriction to nonzero-weight blocks]%
    \label{thm:cpsv_exact_active_reconstruction}
    \lean{MPSTensor.CPSVCanonicalFormData.activeCoordinateMap,
        MPSTensor.CPSVCanonicalFormData.activeCoordinateCoisometry,
        MPSTensor.CPSVCanonicalFormData.activeCoordinateCoisometry_mul_conjTranspose}
    \lean{MPSTensor.CPSVCanonicalFormData.toTensorFromBlocks_active_eq_submatrix,
        MPSTensor.CPSVCanonicalFormData.exact_active_reconstruction}
    \leanok
    \uses{def:cpsv_canonical_form, def:block_diagonal_tensor}
    Suppose the literal CPSV canonical-form data are
    \begin{align}
        A^i        & =V^\dagger\left(\bigoplus_{k=1}^{r}\mu_kA_k^i\right)V,
        \label{eq:cpsv_active_exact_listed}                                 \\
        VV^\dagger & =\Id.
        \label{eq:cpsv_active_exact_listed_coisometry}
    \end{align}
    Put $K_{\mathrm{act}}=\{k:\mu_k\ne0\}$ and
    $D_{\mathrm{act}}=\sum_{k\in K_{\mathrm{act}}}D_k$. There is a matrix
    $U\in\C^{D_{\mathrm{act}}\times D}$ such that the following identities
    hold, the second for every physical index $i$:
    \begin{align}
        UU^\dagger & =\Id,
        \label{eq:cpsv_active_exact_coisometry} \\
        A^i        & =U^\dagger
        \left(\bigoplus_{k\in K_{\mathrm{act}}}\mu_kA_k^i\right)U.
        \label{eq:cpsv_active_exact_reconstruction}
    \end{align}
\end{theorem}

\begin{proof}\leanok
    \uses{def:cpsv_canonical_form, def:block_diagonal_tensor}
    Let $S$ select the flattened coordinates belonging to
    $K_{\mathrm{act}}$. Distinct selected rows are orthonormal, so
    $SS^\dagger=\Id$. The active weighted sum is the principal submatrix
    selected by $S$. Every omitted diagonal block is zero because its weight is
    zero, and therefore
    \begin{align}
        \bigoplus_{k=1}^{r}\mu_kA_k^i
         & =S^\dagger\left(
        \bigoplus_{k\in K_{\mathrm{act}}}\mu_kA_k^i\right)S.
        \label{eq:cpsv_active_exact_selector}
    \end{align}
    Set $U=SV$. Equations~(\ref{eq:cpsv_active_exact_listed}) and
    (\ref{eq:cpsv_active_exact_selector}) give
    (\ref{eq:cpsv_active_exact_reconstruction}), while
    $UU^\dagger=S(VV^\dagger)S^\dagger=\Id$ gives
    (\ref{eq:cpsv_active_exact_coisometry}).
\end{proof}

% Source anchors: CPSV16 lines 217--246, 265--301, and Appendix A,
% lines 1135--1148.
\begin{definition}[Active BNT sector presentation]
    \label{def:cpsv_active_bnt_sector_presentation}
    \lean{MPSTensor.IsActiveCPSVBasisOfNormalTensors}
    \lean{MPSTensor.IsActiveCPSVBasisOfNormalTensors.%
        isCPSVBasisOfNormalTensors}
    \lean{MPSTensor.IsActiveCPSVBasisOfNormalTensors.%
        blocks_not_gaugePhaseEquiv}
    \lean{MPSTensor.IsActiveCPSVBasisOfNormalTensors.of_sameMPV₂Pos}
    \lean{MPSTensor.IsActiveCPSVBasisOfNormalTensors.%
        exists_pos_mpvState_ne_zero}
    \leanok
    \uses{def:sector_weight_data, def:nt_cpsv, def:same_mpv2_pos}
    An \emph{active BNT sector presentation} of a tensor $A$ is a nonempty
    sector decomposition $P$ such that every representative has a positive
    number of copies, every copy has nonzero weight, every representative
    $C_j$ is normal, the representative matrix product vectors are eventually
    linearly independent, and, for every $N>0$,
    \begin{align}
        V^{(N)}(A) & =V^{(N)}(P).
        \label{eq:cpsv_active_bnt_sector_presentation}
    \end{align}
    This records the nonzero canonical summands constructed at
    \cite[lines~217--246]{Cirac2016MPDO_arXiv}: dormant candidates are absent,
    whereas repeated nonzero copies and their weights remain explicit. The
    condition is stronger than the literal BNT definition, which permits the
    zero-coefficient candidate of
    Theorem~\ref{thm:cpsv_bnt_uniqueness_counterexample}.
    % Local fix (dormant BNT candidates): documented in
    % docs/paper-gaps/cpsv16_bnt_uniqueness_zero_coefficient.tex.
\end{definition}

\begin{theorem}[Uniqueness of active BNT presentations]
    \label{thm:cpsv_active_bnt_presentation_uniqueness}
    \lean{MPSTensor.IsActiveCPSVBasisOfNormalTensors.%
        equiv_of_sameMPV₂Pos}
    \leanok
    \uses{def:cpsv_active_bnt_sector_presentation, def:same_mpv2_pos}
    Two active BNT presentations of the same positive-length matrix product
    vector family have equally many representatives. After a permutation,
    corresponding representatives have equal bond dimension and are related
    by an invertible gauge and a unit-modulus phase.

    Only active presentations are compared. The conclusion is false for
    arbitrary literal BNTs by
    Theorem~\ref{thm:cpsv_bnt_uniqueness_counterexample}.
    % Scope restriction (active presentations): documented in
    % docs/paper-gaps/cpsv16_bnt_uniqueness_zero_coefficient.tex.
\end{theorem}

\begin{proof}\leanok
    \uses{thm:cpsv_bnt_characterization, thm:cpsv_bnt_blocks_not_gaugePhaseEquiv, thm:cpsv_normal_mpv_phase_gauge}
    Apply Theorem~\ref{thm:cpsv_bnt_characterization} first to the
    nonzero-weight direct-sum tensor of the first presentation, using the
    representatives of the second as a BNT, and then with the presentations
    interchanged. This gives a matching map in each direction. Eventual linear
    independence excludes gauge-phase-equivalent distinct representatives, so
    both maps are injective. Since the index sets are finite, they have equal
    cardinality, and either matching map is a bijection. Normal-tensor gauge
    rigidity gives the asserted gauge-phase relation for each matched pair.
\end{proof}

% Source anchors: CPSV16 lines 265--301, 317--345, and Appendix C.3,
% lines 1835--1858. Inactive listed blocks remain in the full grouped tensor.
\begin{theorem}[Active representative sector decomposition]
    \label{thm:cpsv_active_representative_sector_decomposition}
    \lean{MPSTensor.CPSVCanonicalFormData.ActiveBNTRefinement.%
        representativeSectorDecomposition}
    \lean{MPSTensor.CPSVCanonicalFormData.ActiveBNTRefinement.%
        representativeSectorDecomposition_basisCount}
    \lean{MPSTensor.CPSVCanonicalFormData.ActiveBNTRefinement.%
        representativeSectorDecomposition_basisDim}
    \lean{MPSTensor.CPSVCanonicalFormData.ActiveBNTRefinement.%
        representativeSectorDecomposition_basis}
    \lean{MPSTensor.CPSVCanonicalFormData.ActiveBNTRefinement.%
        representativeSectorDecomposition_copies}
    \lean{MPSTensor.CPSVCanonicalFormData.ActiveBNTRefinement.%
        representativeSectorDecomposition_weight}
    \lean{MPSTensor.CPSVCanonicalFormData.ActiveBNTRefinement.%
        eventuallyRepresentativeWordTupleSpan}
    \lean{MPSTensor.CPSVCanonicalFormData.ActiveBNTRefinement.%
        groupedTensor_sameMPV₂Pos_representativeSectorDecomposition}
    \lean{MPSTensor.CPSVCanonicalFormData.ActiveBNTRefinement.%
        groupedTensor_isActiveCPSVBasisOfNormalTensors}
    \lean{MPSTensor.CPSVCanonicalFormData.ActiveBNTRefinement.%
        source_sameMPV₂Pos_groupedTensor}
    \lean{MPSTensor.CPSVCanonicalFormData.ActiveBNTRefinement.%
        isActiveCPSVBasisOfNormalTensors}
    \leanok
    \uses{thm:cpsv_canonical_form_active_bnt_refinement, thm:cpsv_bnt_blocked_simultaneous_span, def:cpsv_active_bnt_sector_presentation}
    In the notation of
    Theorem~\ref{thm:cpsv_canonical_form_active_bnt_refinement}, let
    $k(j,q)\in K_{\mathrm{act}}$ be the listed block corresponding to the
    class-copy pair $(j,q)$. The active representatives $C_j$ define a sector
    decomposition $P_{\mathrm{act}}$ with copy weights
    \begin{align}
        \lambda_{j,q}
         & =\nu_{k(j,q)}\zeta_{k(j,q)}.
        \label{eq:cpsv_active_representative_weight}
    \end{align}
    Every copy weight in $P_{\mathrm{act}}$ is nonzero. The representatives have simultaneous
    word-tuple separation at every sufficiently large length.

    The sector decomposition contains only the active class-copy pairs. The
    full grouped tensor $T_{\mathrm{grp}}$ of
    Theorem~\ref{thm:cpsv_canonical_form_active_bnt_refinement} still contains
    every inactive zero-weight listed coordinate. For every positive length
    $N$, the grouped and active-sector tensors generate the same matrix product
    vector. The original tensor has the same positive-length matrix product
    vectors as both:
    \begin{align}
        V^{(N)}(T_{\mathrm{grp}})
         & =V^{(N)}(P_{\mathrm{act}}).
        \label{eq:cpsv_active_representative_positive_mpv}
        \\
        V^{(N)}(A)
         & =V^{(N)}(P_{\mathrm{act}}).
        \label{eq:cpsv_active_representative_source_mpv}
    \end{align}
    Hence $P_{\mathrm{act}}$ is an active BNT sector presentation of $A$
    whenever at least one active listed block occurs.
    No equality at length zero is asserted, since deleting inactive bond
    coordinates changes the empty-word coefficient.
\end{theorem}

\begin{proof}\leanok
    \uses{thm:cpsv_canonical_form_active_bnt_refinement, thm:cpsv_bnt_blocked_simultaneous_span, thm:coisometry_reconstruction_same_mpv_pos, thm:gauge_invariance, lem:paper_decomp_formula}
    The active class-copy equivalence reindexes the nonzero listed blocks.
    (\ref{eq:cpsv_active_representative_weight}) combines the original
    coefficient with the unit phase relating its block to $C_j$; both factors
    are nonzero. The basis-of-normal-tensors condition and
    Theorem~\ref{thm:cpsv_bnt_blocked_simultaneous_span} give eventual
    simultaneous word-tuple separation for the representatives. At positive
    length, every inactive contribution contains a positive power of its zero
    coefficient and vanishes. Reindexing the remaining finite sum by $(j,q)$
    gives (\ref{eq:cpsv_active_representative_positive_mpv}).
    The original coisometric reconstruction and simultaneous block gauge
    preserve every positive-length closed-chain coefficient, yielding
    (\ref{eq:cpsv_active_representative_source_mpv}).
\end{proof}

\begin{definition}[BNT canonical form]%
    \label{def:sector_bnt_cf}
    \lean{MPSTensor.IsBNTCanonicalForm}
    \leanok
    \uses{def:sector_weight_data}
    Given a sector decomposition $P$, these conditions define the normalized,
    left-canonical, aperiodic BNT surface used below. This is stronger than the
    abstract BNT of Definition~\ref{def:bnt} and the literal coefficient-form
    BNT of Definition~\ref{def:bnt_cpsv}: besides spanning and eventual linear
    independence, it fixes irreducible left-canonical representatives, their
    self-overlap normalization, gauge-phase separation, and the canonical-form
    weight normalization. It imposes no equal-modulus factorization or strict
    ordering on the raw weights. Explicitly, each basis block has positive bond
    dimension, is irreducible, and is left-canonical; the normalized self-overlap of every basis block
    converges to $1$; the basis MPV family is eventually linearly
    independent; distinct same-dimension basis blocks are not gauge-phase
    equivalent after identifying equal bond dimensions; and the
    \cite[Section~2.3]{Cirac2016MPDO_arXiv} canonical-form normalization holds,
    \begin{align}
        \forall (j,q),\quad       & |\mu_{j,q}| \le 1,
        \notag                                           \\
        \exists\, (j_*,q_*),\quad & |\mu_{j_*,q_*}| = 1.
        \notag
    \end{align}
    The MPV expansion takes the raw two-layer form
    \begin{align}
        V^{(N)}(P)_\sigma
         & =
        \sum_{j=1}^{g}
        \left(\sum_{q=1}^{r_j} \mu_{j,q}^N\right)
        V^{(N)}(P_j)_\sigma,
        \label{eq:bnt_sector_raw_expansion}
    \end{align}
    % Source: Cirac2016MPDO_arXiv \S~II.C lines 271--301; Cirac2021Matrix
    % Definition~4.2, Proposition~4.3, lines 1864--1884.
    matching \cite[Section~2.3]{Cirac2016MPDO_arXiv} and
    \cite[Definition~IV.2]{Cirac2021Matrix}.
    The canonical-form normalization admits every example of the source paper
    (in particular $C \oplus D$, $C \oplus (-C)$,
    $C \oplus (1/2) C$, and $C \oplus (-C) \oplus (1/2) C$) and is
    not derivable from the per-block normality hypotheses alone.
\end{definition}

\begin{definition}[Single-sector decomposition]
    \label{def:single_sector_decomposition}
    \lean{MPSTensor.singleSectorDecomposition}
    \leanok
    \uses{def:sector_weight_data}
    For a tensor $A$, its single-sector decomposition has $A$ as its only
    basis tensor, with one copy of weight $1$.
\end{definition}

\begin{theorem}[Single-sector BNT canonical form]
    \label{thm:single_sector_decomposition_bnt}
    \lean{MPSTensor.isBNTCanonicalForm_singleSectorDecomposition}
    \leanok
    \uses{def:sector_bnt_cf, def:single_sector_decomposition}
    Let $A$ have positive bond dimension, be irreducible and left-canonical,
    and satisfy $O_{AA}(N)\longrightarrow1$. Then its single-sector
    decomposition is a BNT canonical form.
\end{theorem}

\begin{proof}\leanok
    The identity
    \begin{align}
        O_{AA}(N)
         & =\left\lVert V^{(N)}(A)\right\rVert^2
        \longrightarrow 1\ne0.
    \end{align}
    shows that $V^{(N)}(A)$ is nonzero for all sufficiently large $N$.
    A family consisting of one nonzero vector is linearly independent. All
    remaining conditions are immediate from the hypotheses and the unit
    weight.
\end{proof}

\begin{theorem}[Normal tensors have single-sector BNT presentations]
    \label{thm:normal_single_sector_bnt_gauge}
    \lean{MPSTensor.IsNormalTensor.exists_gaugeEquiv_singleSectorDecomposition}
    \leanok
    \uses{def:nt_cpsv, def:gauge_equiv, def:sector_bnt_cf, def:single_sector_decomposition, def:same_mpv2_pos}
    Every normal tensor is gauge equivalent to a tensor whose single-sector
    decomposition is a BNT canonical form. The original tensor and the total
    tensor of this decomposition generate the same matrix product vectors at
    every positive length.
\end{theorem}

\begin{proof}\leanok
    \uses{thm:single_sector_decomposition_bnt, thm:cpsv_normal_tp_gauge, thm:peripheral_primitive_irreducible_overlap_one, thm:gauge_invariance, lem:paper_decomp_formula}
    Choose the pure trace-preserving Perron gauge. Its tensor is irreducible
    and left-canonical, and primitivity gives self-overlap converging to one.
    The preceding theorem supplies the BNT canonical form. If $P$ denotes
    the total tensor of its single-sector decomposition, then gauge invariance
    and the sector-decomposition formula give, for every $N>0$ and word
    $\sigma$,
    \begin{align}
        V^{(N)}(A)_\sigma
         & =V^{(N)}(B)_\sigma
        =1^N V^{(N)}(B)_\sigma
        =V^{(N)}(P)_\sigma.
    \end{align}
\end{proof}

\begin{lemma}[Unit-modulus weight witness from the canonical-form normalization]%
    \label{lem:sector_bnt_weight_unit_exists_of_struct}
    \lean{MPSTensor.IsBNTCanonicalForm.weight_unit_exists_of_struct}
    \leanok
    \uses{def:sector_bnt_cf}
    Under the BNT canonical form, the
    \cite[Section~2.3]{Cirac2016MPDO_arXiv} normalization provides indices
    $(j_*,q_*)$ with $|\mu_{j_*,q_*}|=1$.
\end{lemma}

\begin{proof}\leanok
    This is the unit-modulus existential in the canonical-form normalization.
\end{proof}

\begin{corollary}[Basis blocks of a BNT canonical form are normal]
    \label{cor:sector_bnt_cf_basis_normal}
    \lean{MPSTensor.IsBNTCanonicalForm.basis_isNormal}
    \leanok
    \uses{def:sector_bnt_cf, thm:bnt_normality_from_self_overlap}
    If $P$ is in BNT canonical form, then every basis block $P_j$ is normal
    in the algebraic sense of Definition~\ref{def:normal}.
\end{corollary}

\begin{proof}\leanok
    \uses{thm:bnt_normality_from_self_overlap}
    Each basis block has positive bond dimension, is irreducible and
    left-canonical, and has normalized self-overlap converging to $1$.
    Theorem~\ref{thm:bnt_normality_from_self_overlap} applies to each block.
\end{proof}

\begin{theorem}[A BNT canonical form determines an algebraic BNT]
    \label{thm:sector_bnt_cf_is_bnt}
    \lean{MPSTensor.IsBNTCanonicalForm.isBNT}
    \leanok
    \uses{def:bnt, def:sector_bnt_cf, cor:sector_bnt_cf_basis_normal}
    If $P$ is in BNT canonical form, then its basis blocks form an algebraic
    basis of normal tensors for the total tensor associated with $P$. Only
    this implication from canonical form to algebraic BNT is asserted.
\end{theorem}

\begin{proof}\leanok
    \uses{cor:sector_bnt_cf_basis_normal, def:sector_weight_data}
    The preceding corollary gives normality of every basis block. The sector
    coefficient formula (\ref{eq:bnt_sector_decomp}) expresses the total MPV
    as a linear combination of the basis MPVs at each positive length, and the
    BNT sector data give their eventual linear independence.
\end{proof}

\begin{theorem}[Positive blocking preserves BNT canonical form]
    \label{thm:sector_bnt_cf_block_tensor}
    \lean{MPSTensor.SectorDecomposition.IsBNTCanonicalForm.blockTensor}
    \leanok
    \uses{def:sector_decomposition_block_tensor, def:sector_bnt_cf, thm:bnt_primitivity_normality_from_self_overlap}
    Let $P$ be a BNT canonical form and let $p\geq1$. Then the blocked
    sector decomposition $P^{[p]}$ is again a BNT canonical form.
\end{theorem}

\begin{proof}
    \leanok
    \uses{thm:sector_decomposition_block_tensor_mpv, thm:bnt_primitivity_normality_from_self_overlap}
    For each representative $A_j$, normalized self-overlap gives peripheral
    period one, hence a primitive transfer map. Positive blocking therefore
    preserves irreducibility and left-canonical normalization. Moreover,
    $O_{A_j^{[p]}A_j^{[p]}}(N)=O_{A_jA_j}(Np)\longrightarrow1$.
    Let $F_N$ be the linear isomorphism induced by concatenating blocked words,
    from the blocked configuration space at length $N$ to the original one at
    length $Np$. For every representative,
    $F_N\ket{V^{(N)}(A_j^{[p]})}=\ket{V^{(Np)}(A_j)}$.
    Hence eventual linear independence is preserved. If two blocked
    representatives were gauge-phase equivalent, then for some nonzero
    $\zeta$ and every $N$,
    $\ket{V^{(N)}(A_k^{[p]})}
        =\zeta^N\ket{V^{(N)}(A_j^{[p]})}$.
    Writing these two vectors as $f_k$ and $f_j$ gives
    $1\cdot f_k=\zeta^N\cdot f_j$ in an eventually linearly independent
    family, and therefore $k=j$. Finally,
    $|\mu_{j,q}^p|=|\mu_{j,q}|^p\leq1$, and a unit-modulus weight remains of
    unit modulus after taking its $p$-th power.
\end{proof}

\begin{theorem}[Common exact quantum-Wielandt length for BNT representatives]
    \label{thm:bnt_representatives_exact_total_dimension_pow_four}
    \lean{MPSTensor.IsBNTCanonicalForm.basis_isNBlkInjective_totalDim_pow_four}
    \leanok
    \uses{def:sector_bnt_cf, thm:sector_basis_dimension_le_total_dimension, thm:bnt_normality_from_self_overlap, thm:exact_pow_four_block_injectivity_normal_left_canonical, thm:wordspan_blk_inj_succ}
    Let $P$ be a BNT canonical form of total bond dimension $D$, with
    representatives $A_j$ of bond dimensions $D_j$. For every
    $0\leq j<g$, the representative $A_j$ is $D^4$-block injective:
    $S_{D^4}(A_j)=\MN{D_j}$.
\end{theorem}

\begin{proof}
    \leanok
    \uses{def:sector_bnt_cf, thm:sector_basis_dimension_le_total_dimension, thm:bnt_normality_from_self_overlap, thm:exact_pow_four_block_injectivity_normal_left_canonical, thm:wordspan_blk_inj_succ}
    The BNT hypotheses make $A_j$ irreducible and left-canonical, with
    $O_{A_jA_j}(N)\to1$. The preceding theorem therefore makes $A_j$ normal,
    and quantum Wielandt gives $S_{D_j^4}(A_j)=\MN{D_j}$.
    By Theorem~\ref{thm:sector_basis_dimension_le_total_dimension},
    $D_j^4\leq D^4$. Write $D^4=D_j^4+k$ and apply positive-length
    propagation $k$ times.
\end{proof}

\begin{theorem}[CPSV simultaneous block-injectivity bound]
    \label{thm:sharp_bnt_block_separation}
    \lean{MPSTensor.IsBNTCanonicalForm.exists_basis_wordTupleSpanTop_le_three_totalDim_pow_five}
    \lean{MPSTensor.three_mul_pred_mul_pow_four_add_one_le_three_mul_pow_five}
    \leanok
    \uses{def:sector_bnt_cf, def:blockwise_product_word_span, thm:bnt_representatives_exact_total_dimension_pow_four, lem:bnt_direct_sum_block_separation_word_span}
    Let $P$ be a BNT canonical form with total bond dimension $D$. There is a
    positive integer $N\leq3D^5$ such that its representatives have the
    simultaneous fixed-length span
    \begin{align}
        \spn\{(A_1^\omega,\ldots,A_g^\omega):|\omega|=N\}
         & =\prod_{j=1}^{g}\MN{D_j}.
        \notag
    \end{align}
    This is the block-injectivity estimate in
    \cite[lines~340--345]{Cirac2016MPDO_arXiv}.
\end{theorem}

\begin{proof}
    \leanok
    \uses{thm:bnt_representatives_exact_total_dimension_pow_four, lem:bnt_direct_sum_block_separation_word_span, thm:sector_basis_dimension_le_total_dimension}
    Every representative is block-injective at $D^4$, and homogeneous-span
    propagation gives the same conclusion at $D^4+1$ and $3(D^4+1)$. If
    $g=1$, take $N=D^4$. If $g\geq2$, apply
    Lemma~\ref{lem:bnt_direct_sum_block_separation_word_span} with $L_0=D^4$ and take
    \begin{align}
        N & =3(g-1)(D^4+1).
        \notag
    \end{align}
    The positive bond dimension of every representative gives $g\leq D$.
    Consequently $N\leq3D^5$; the same estimate also covers the one-block
    case. The unit-modulus weight required by the BNT normalization ensures
    $D>0$.
\end{proof}

The matching arguments below read off three facts from a BNT canonical form: the
raw coefficient identity, the cross-overlap decay between distinct basis blocks,
and the combined-family linear independence.

\begin{lemma}[Raw two-layer sector coefficient identity]%
    \label{lem:sector_bnt_coeff_eq_sum_weight_pow}
    \lean{MPSTensor.SectorDecomposition.coeff_eq_sum_weight_pow}
    \leanok
    \uses{def:sector_weight_data}
    For every sector decomposition $P$, length $N$, and basis index $j$,
    $c_N^{(j)} = \sum_{q=1}^{r_j} \mu_{j,q}^N$.
    % Source: two-layer BNT and coefficient-expansion displays,
    % Cirac2016MPDO_arXiv lines 287--301; Cirac2021Matrix lines 1864--1884.
    This is \cite[Section~2.3]{Cirac2016MPDO_arXiv}.
\end{lemma}

\begin{proof}\leanok
    The claim follows from the definition of the sector coefficient.
\end{proof}

\begin{lemma}[Cross-overlap decay between distinct basis blocks]%
    \label{lem:sector_bnt_cross_overlap_basis_tendsto_zero}
    \lean{MPSTensor.IsBNTCanonicalForm.cross_overlap_basis_tendsto_zero}
    \leanok
    \uses{def:mpv_overlap, def:sector_bnt_cf}
    If $P$ is a BNT canonical form and $j \neq k$, then
    $O_{P_j P_k}(N) \to 0$.
    % Source: normal-tensor overlap dichotomy, Cirac2016MPDO_arXiv lines
    % 1080--1091, with the gauge-phase distinctness clause at lines 264--279.
\end{lemma}

\begin{proof}\leanok
    Case on the bond dimensions. If they differ, the dimension-mismatch
    overlap-decay theorem applies. If they agree, the basis-distinctness clause
    rules out gauge-phase equivalence, so the equal-dimension overlap-decay
    theorem applies.
\end{proof}

\begin{lemma}[Combined-family eventual linear independence]%
    \label{lem:sector_bnt_combined_family_eventually_li}
    \lean{MPSTensor.IsBNTCanonicalForm.combined_family_eventually_li}
    \leanok
    \uses{lem:bnt_two_family_overlap_orthonormal_li, lem:sector_bnt_cross_overlap_basis_tendsto_zero}
    Let $P$ and $Q$ each be a BNT canonical form, and suppose every mixed overlap
    $O_{P_j Q_k}(N)$ tends to $0$. Then the union family
    \begin{align}
        \bigl\{\ket{V^{(N)}(P_j)}\bigr\}_{j=1}^{g_P}
        \cup
        \bigl\{\ket{V^{(N)}(Q_k)}\bigr\}_{k=1}^{g_Q}
        \notag
    \end{align}
    is linearly independent for all sufficiently large $N$.
    % Source: corollary at Cirac2016MPDO_arXiv lines 1121--1132; the BNT
    % linear-independence input is recorded at Cirac2021Matrix line 1850.
    This is \cite[Appendix~A]{Cirac2016MPDO_arXiv}.
\end{lemma}

\begin{proof}\leanok
    Feed the per-family normalized self-overlaps, the within-family
    cross-overlap decay
    (Lemma~\ref{lem:sector_bnt_cross_overlap_basis_tendsto_zero}), and the
    inter-family decay hypothesis into the two-family overlap-orthonormality
    criterion (Lemma~\ref{lem:bnt_two_family_overlap_orthonormal_li}).
\end{proof}

\begin{example}[$C \oplus (1/2)C$, unequal-modulus admissibility]%
    \label{ex:sector_bnt_halvedDecomp}
    \lean{MPSTensor.SectorBNT.Examples.halvedDecomp}
    \leanok
    \uses{def:sector_weight_data}
    For any tensor $C$, this construction forms a single-sector, two-copy
    decomposition with basis tensor $C$ and raw weights $(1,1/2)$. Thus the
    sector data themselves allow unequal-modulus copies. If $C$ additionally
    has positive bond dimension, is irreducible and left-canonical, has
    normalized nonzero self-overlap, and satisfies the BNT independence clause,
    these data satisfy Definition~\ref{def:sector_bnt_cf}; those hypotheses are
    not consequences of the construction alone.
\end{example}

\subsection{Single-copy sector decompositions}

A weighted block sum is the one-copy case of the sector decomposition used
below. Equal or repeated sectors require the full coefficient formalism.

\subsection{Prepared-block construction of BNT canonical forms}%
\label{subsec:paperbnt_supplier}

The construction below separates the after-blocking preparation of suitable
blocks from the normalized-weight BNT construction: the prepared blocks are
trace-preserving, primitive, irreducible, and injective; normalized weights then
quotient gauge-phase-equivalent blocks into sectors to produce the canonical
form of Definition~\ref{def:sector_bnt_cf}.

\begin{theorem}[Prepared BNT blocks after blocking]%
    \label{thm:prepared_bnt_blocks_after_blocking}
    \lean{MPSTensor.exists_prepared_BNT_blocks_afterBlocking_pos}
    \leanok
    \uses{thm:unconditional_common_primitive_irreducible_blocks, thm:common_injective_blocking_tp_primitive_irreducible_family}
    Let $A$ be an MPS tensor. Then there is a blocking length $p\ge 1$, nonzero
    weights $\mu_k$, and blocks $B_k$ over the $p$-blocked alphabet such that:
    \begin{enumerate}
        \item each block has positive bond dimension;
        \item each block is trace-preserving, primitive, tensor-irreducible, and
              one-site injective;
        \item $A^{[p]}$ and $\bigoplus_k \mu_k B_k$ generate the same MPV family
              at every positive length.
    \end{enumerate}
    No weight-modulus normalization is asserted at this stage.
\end{theorem}

\begin{proof}\leanok
    Apply the common primitive irreducible block decomposition to $(A,A)$,
    giving a positive blocking length and a nonzero weighted family of
    trace-preserving primitive irreducible blocks with the positive-length MPV
    identity. Theorem~\ref{thm:common_injective_blocking_tp_primitive_irreducible_family}
    supplies one common further blocking making every block one-site injective
    while preserving the other hypotheses; finally identify the iterated blocked
    alphabet with the direct one.
\end{proof}

With the prepared blocks in hand, normalized weights produce a BNT canonical
form by quotienting gauge-phase-equivalent blocks into sectors.

\begin{theorem}[BNT canonical form of the phase-class quotient]%
    \label{thm:paperbnt_collapsed_prepared}
    \lean{MPSTensor.isBNTCanonicalForm_collapsedBntSectorDecomp_of_tp_primitive_irr_blocks}
    \leanok
    \uses{def:mpv_phase_class_data, def:sector_bnt_cf}
    Let $(B_k)_{k=1}^r$ have positive bond dimensions and suppose that every
    $B_k$ is trace-preserving, primitive, and irreducible. Assume
    $\mu_k\ne0$ and $|\mu_k|\le1$ for every $k$, and
    $|\mu_{k_*}|=1$ for some $k_*$. Partition the blocks into
    matrix-product-vector phase classes, choose one
    representative from each class, and retain every class member as a copy of
    its representative with the corresponding phase-adjusted weight. The
    resulting sector decomposition is in BNT canonical form.
\end{theorem}

\begin{proof}\leanok
    \uses{def:mpv_phase_class_data, def:sector_bnt_cf}
    The chosen representatives inherit irreducibility and trace preservation,
    while primitivity gives normalized self-overlap. Distinct phase classes
    give gauge-phase-distinct representatives. The collapsed phase-class
    construction supplies the eventual linear independence and retains every
    original block as one copy. Its phase multipliers have unit modulus, so the
    adjusted weights preserve the assumed bounds and unit witness. These are
    precisely the fields of Definition~\ref{def:sector_bnt_cf}.
\end{proof}

\begin{theorem}[BNT from normal blocks with normalized weights]%
    \label{thm:paperbnt_supplier_prepared}
    \lean{MPSTensor.exists_isBNTCanonicalForm_of_tp_primitive_irr_blocks}
    \leanok
    Let $r \in \N$, and for each $k = 1,\ldots,r$ let $D_k \in \N$,
    $\mu_k \in \C$, and let $B_k$ be an MPS tensor of physical dimension $d$ and
    bond dimension $D_k$. Assume:
    \begin{enumerate}
        \item \textbf{positive bond dimensions:} $D_k > 0$ for all $k$,
        \item \textbf{TP normalization:} $\sum_i (B_k^i)^\dagger B_k^i = \Id$,
        \item \textbf{primitivity:} $\E_{B_k}$ is primitive for each $k$,
        \item \textbf{irreducibility:} each $B_k$ is irreducible,
        \item \textbf{nonzero weights:} $\mu_k \neq 0$ for all $k$,
        \item \textbf{bounded weights:} $|\mu_k| \le 1$ for all $k$,
        \item \textbf{global unit witness:} $|\mu_k| = 1$ for some $k$.
    \end{enumerate}
    Then there exists a sector decomposition $P$ such that $P$ and
    $\bigoplus_{k=1}^{r} \mu_k B_k$ generate the same MPV family at every
    positive length,
    and $P$ satisfies the basis-of-normal-tensors conditions
    (Definition~\ref{def:bnt_cpsv}) together with the multiplicity, span, and
    weight-bound conditions of Definition~\ref{def:sector_bnt_cf}.
\end{theorem}

\begin{proof}\leanok
    The phase-class quotient groups gauge-phase-equivalent blocks into sectors
    with unit-modulus per-copy scalars. Under the TP, primitive, and irreducible
    hypotheses the phase-class construction applies, and for distinct
    representatives $B_j$, $B_{j'}$ ($j \neq j'$) the cross-overlap decays,
    \begin{align}
        \lim_{N\to\infty} O_{B_j B_{j'}}(N) & =0,
        \notag                                    \\
        \lim_{N\to\infty} O_{B_j B_j}(N)    & =1,
        \notag
    \end{align}
    which gives the eventual linear-independence condition with no injectivity
    required. The resulting sector decomposition inherits the span and
    multiplicity conditions, and the bounded-weight and global-unit hypotheses
    give the weight-norm conditions.
\end{proof}

The BNT construction preserves total bond dimension because CF summands assigned
to the same BNT representative have equal bond dimension. Indeed, suppose
trace-preserving primitive irreducible NTs $X$ and $Y$ of positive bond dimension
have MPV families related by
$\ket{V^{(N)}(Y)}=\zeta^N\ket{V^{(N)}(X)}$ with $\zeta\neq0$. The normalized
self-overlaps give $|O_{XX}(N)|\to 1$. The phase relation gives the
cross-overlap identity
\begin{align}
    O_{YX}(N)   & =\zeta^N O_{XX}(N),
    \notag                                \\
    |O_{YX}(N)| & =|\zeta|^N |O_{XX}(N)|.
    \notag
\end{align}
Since the same normalization forces $|\zeta|=1$, this gives
$|O_{YX}(N)|\to 1$.
If their bond dimensions differed, the rectangular overlap-decay theorem would
give $O_{YX}(N)\to 0$, a contradiction.

\begin{theorem}[Gauge realization of prepared phase-equivalent NTs]%
    \label{thm:prepared_phase_equiv_gauge}
    \lean{MPSTensor.gaugePhaseEquiv_of_MPVBlockPhaseEquiv_of_tp_primitive_irr}
    \leanok
    \uses{def:mpv_block_phase_equiv, def:gauge_phase_equiv, def:irreducible_tensor}
    Let $X$ and $Y$ be trace-preserving, primitive, irreducible NTs of positive
    bond dimensions. If their matrix-product-vector families differ by a
    nonzero scalar power, then their bond dimensions are equal and, after
    identifying the bond spaces, $X$ and $Y$ are gauge-phase equivalent.
    Neither positivity of the scalar nor unitarity of the gauge is asserted.
\end{theorem}

\begin{proof}\leanok
    \uses{lem:prepared_phase_equiv_dim_eq, thm:not_gauge_phase_eventually_not_proportional, thm:peripheral_primitive_irreducible_overlap_one}
    Lemma~\ref{lem:prepared_phase_equiv_dim_eq} identifies the bond dimensions.
    If the resulting same-dimensional blocks were not gauge-phase equivalent,
    Theorem~\ref{thm:not_gauge_phase_eventually_not_proportional} would give a
    positive length at which their matrix product vectors are not proportional.
    This contradicts the assumed scalar-power relation. The normalized
    self-overlaps required there follow from
    Theorem~\ref{thm:peripheral_primitive_irreducible_overlap_one}.
\end{proof}

\begin{lemma}[Gauge realization with the fixed phase scalar]%
    \label{lem:prepared_phase_equiv_fixed_scalar_gauge}
    \lean{MPSTensor.exists_gauge_choose_MPVBlockPhaseEquiv_of_tp_primitive_irr}
    \leanok
    \uses{def:mpv_block_phase_equiv, def:gauge_phase_equiv, def:irreducible_tensor}
    Let $X$ and $Y$ be trace-preserving, primitive, irreducible NTs of positive
    bond dimensions. Fix $\zeta\in\C\setminus\{0\}$ and suppose that, for every
    positive length $N$
    and every physical word $\sigma$,
    \begin{align}
        V^{(N)}(Y)_\sigma
         & =\zeta^N V^{(N)}(X)_\sigma.
        \label{eq:prepared_phase_equiv_fixed_scalar}
    \end{align}
    Then, after identifying the bond spaces, there exists
    $G\in\GL_{D_Y}(\C)$ such that, for every physical index $i$,
    \begin{align}
        Y^i & =\zeta G X^i G^{-1}.
        \label{eq:prepared_phase_equiv_fixed_scalar_gauge}
    \end{align}
\end{lemma}

\begin{proof}\leanok
    \uses{thm:prepared_phase_equiv_gauge, thm:peripheral_primitive_irreducible_overlap_one, thm:gauge_phase_mpv_scaling}
    Theorem~\ref{thm:prepared_phase_equiv_gauge} gives an invertible $G$ and a
    nonzero scalar $\eta$ such that $Y^i=\eta G X^iG^{-1}$. By
    Theorem~\ref{thm:gauge_phase_mpv_scaling},
    $V^{(N)}(Y)_\sigma=\eta^N V^{(N)}(X)_\sigma$.
    Theorem~\ref{thm:peripheral_primitive_irreducible_overlap_one} gives
    consecutive positive lengths $M$ and $M+1$ with nonzero coefficients of
    $V^{(M)}(X)$ and $V^{(M+1)}(X)$. Comparing the two scalar-power identities
    gives $\zeta^M=\eta^M$ and $\zeta^{M+1}=\eta^{M+1}$. Hence
    $\zeta=\eta$, which proves
    (\ref{eq:prepared_phase_equiv_fixed_scalar_gauge}).
\end{proof}

\begin{theorem}[Dimension-preserving BNT form of prepared blocks]%
    \label{thm:paperbnt_supplier_prepared_total_dim}
    \lean{MPSTensor.exists_isBNTCanonicalForm_of_tp_primitive_irr_blocks_and_totalDim}
    \leanok
    \uses{thm:paperbnt_collapsed_prepared, thm:prepared_phase_equiv_gauge, def:sector_weight_data}
    Under the hypotheses of
    Theorem~\ref{thm:paperbnt_collapsed_prepared}, there is a BNT canonical-form
    sector decomposition $P$ such that
    \begin{align}
        \mc{V}^+(P)
         & =\mc{V}^+\!\left(\bigoplus_{k=1}^r\mu_kB_k\right),
        \notag                                                \\
        D_{\mathrm{tot}}(P)
         & =\sum_{k=1}^r D_k.
        \label{eq:paperbnt_prepared_total_dim}
    \end{align}
\end{theorem}

\begin{proof}\leanok
    \uses{thm:paperbnt_collapsed_prepared, thm:prepared_phase_equiv_gauge}
    Use the phase-class sector decomposition from
    Theorem~\ref{thm:paperbnt_collapsed_prepared}. Its flattened copies are in
    bijection with the original blocks, and phase-equivalent blocks have equal
    bond dimension by Theorem~\ref{thm:prepared_phase_equiv_gauge}. Summing the
    copy dimensions over this bijection gives
    (\ref{eq:paperbnt_prepared_total_dim}).
\end{proof}

\begin{theorem}[Exact reindexing of heterogeneous weighted block sums]%
    \label{thm:heterogeneous_block_sum_reindex}
    \lean{MPSTensor.blockDimEquiv,
        MPSTensor.toTensorFromBlocks_eq_reindex_of_equiv}
    \leanok
    \uses{def:block_diagonal_tensor}
    Let $e:\{1,\ldots,r\}\simeq\{1,\ldots,r'\}$ satisfy
    $D_k=D'_{e(k)}$. Let $Q_k$ be the coordinate permutation whose rows are
    indexed by the primed block and whose columns are indexed by the unprimed
    block. Suppose that, for every block and physical index,
    \begin{align}
        \mu_kA_k^i
         & =Q_k^\dagger\left(\mu'_{e(k)}B_{e(k)}^i\right)Q_k.
        \label{eq:heterogeneous_block_local}
    \end{align}
    The induced equivalence $E$ of flattened dependent block coordinates has a
    permutation matrix $Q_E$, with rows indexed by the primed coordinates and
    columns indexed by the unprimed coordinates. For every physical index $i$,
    \begin{align}
        \bigoplus_{k=1}^{r}\mu_kA_k^i
         & =Q_E^\dagger\left(\bigoplus_{\ell=1}^{r'}
        \mu'_{\ell}B_{\ell}^i\right)Q_E.
        \label{eq:heterogeneous_block_global}
    \end{align}
\end{theorem}

\begin{proof}\leanok
    \uses{def:block_diagonal_tensor}
    Evaluate both sides on two flattened block coordinates. When their block
    labels agree, (\ref{eq:heterogeneous_block_local}) gives the matrix entry.
    When the labels differ, both block-diagonal entries vanish. This proves
    (\ref{eq:heterogeneous_block_global}).
\end{proof}

\begin{theorem}[Exact BNT reconstruction of prepared normal blocks]%
    \label{thm:paperbnt_supplier_prepared_exact}
    \lean{MPSTensor.exists_isBNTCanonicalForm_exact_of_tp_primitive_irr_blocks}
    \leanok
    \uses{thm:paperbnt_collapsed_prepared, def:block_diagonal_tensor, def:sector_weight_data}
    Under the hypotheses of
    Theorem~\ref{thm:paperbnt_collapsed_prepared}, there are a BNT
    canonical-form sector decomposition $P$, an equivalence $E$ of flattened
    bond coordinates, and $X\in\GL_{D_{\mathrm{tot}}(P)}(\C)$. Write $Q_E$ for
    the permutation matrix whose rows are indexed by the coordinates of $P$ and
    whose columns are indexed by the original block coordinates. The first
    identity below holds, and the second holds for every physical index $i$:
    \begin{align}
        D_{\mathrm{tot}}(P) & =\sum_{k=1}^{r}D_k,
        \label{eq:prepared_exact_total_dim}                           \\
        \bigoplus_{k=1}^{r}\mu_kB_k^i
                            & =Q_E^\dagger\left(XP^iX^{-1}\right)Q_E.
        \label{eq:prepared_exact_reconstruction}
    \end{align}
\end{theorem}

\begin{proof}\leanok
    \uses{def:mpv_phase_class_data, thm:paperbnt_collapsed_prepared, lem:prepared_collapsed_bnt_sector_decomp_total_dim, lem:prepared_phase_equiv_fixed_scalar_gauge, lem:tensor_from_blocks_globalGauge_conj, thm:heterogeneous_block_sum_reindex, def:global_gauge_of_blocks}
    Partition the prepared blocks into matrix-product-vector phase classes. For
    the copy $q$ in class $j$, retain the precise class scalar $\zeta_{j,q}$
    chosen in the weight $\zeta_{j,q}\mu_{j,q}$. The normal-tensor phase
    theorem supplies an invertible $G_{j,q}$ with
    \begin{align}
        B_{j,q}^i & =\zeta_{j,q}G_{j,q}C_j^iG_{j,q}^{-1}.
        \label{eq:prepared_exact_copy_gauge}
    \end{align}
    Thus the original weighted block equals the corresponding sector copy after
    conjugation, since its sector weight is
    $\zeta_{j,q}\mu_{j,q}$. Assemble the $G_{j,q}$ into a block-diagonal
    matrix $X$ and apply
    Theorem~\ref{thm:heterogeneous_block_sum_reindex}. This gives
    (\ref{eq:prepared_exact_reconstruction}); summing the matched copy
    dimensions gives (\ref{eq:prepared_exact_total_dim}).
\end{proof}

\begin{lemma}[Gauge equivalence of direct sums from per-block gauge equivalence]%
    \label{lem:gauge_equiv_from_block_gauge}
    \lean{MPSTensor.gaugeEquiv_toTensorFromBlocks_of_blockGauge}\leanok
    \uses{lem:tensor_from_blocks_globalGauge_conj, def:gauge_equiv}
    If each pair $(A_k, B_k)$ is gauge equivalent, then the weighted
    block-diagonal tensors $\bigoplus_k \mu_k A_k$ and
    $\bigoplus_k \mu_k B_k$ are gauge equivalent.
\end{lemma}

\begin{proof}\leanok
    Choose per-block gauge matrices via the hypothesis; then apply
    Lemma~\ref{lem:tensor_from_blocks_globalGauge_conj}.
\end{proof}

\begin{theorem}[Positive-length MPVs under coisometric reconstruction]%
    \label{thm:coisometry_reconstruction_same_mpv_pos}
    \lean{MPSTensor.sameMPV₂Pos_of_coisometry_reconstruction}
    \leanok
    \uses{def:same_mpv2_pos}
    Let $A$ and $B$ be MPS tensors, possibly with different bond dimensions.
    Suppose that $U$ satisfies $UU^\dagger=\Id$ and
    $A^i=U^\dagger B^iU$ for every physical index $i$. Then $A$ and $B$
    generate the same matrix product vectors at every positive length.
\end{theorem}

\begin{proof}\leanok
    \uses{def:mpv}
    Multiply the reconstruction identities around a positive-length cycle.
    Each adjacent pair $UU^\dagger$ cancels, and cyclicity of the trace removes
    the remaining outer factors.
\end{proof}

\begin{theorem}[Active-sector BNT form of normalized CPSV canonical data]%
    \label{thm:cpsv_active_sector_bnt}
    \lean{MPSTensor.CPSVCanonicalFormData.exists_active_isBNTCanonicalForm}
    \leanok
    \uses{def:cpsv_canonical_form, def:sector_bnt_cf, def:same_mpv2_pos}
    Let $A$ have CPSV canonical-form data satisfying the separate normalization
    condition~(\ref{eq:cpsv_weight_normalization}), and assume $A\ne0$. Put
    $K_{\mathrm{act}}=\{k:\mu_k\ne0\}$. Then there is a BNT canonical-form
    sector decomposition $P$ such that
    \begin{align}
        \mc{V}^+(A) & =\mc{V}^+(P),
        \notag                                          \\
        D_{\mathrm{tot}}(P)
                    & =\sum_{k\in K_{\mathrm{act}}}D_k.
        \label{eq:cpsv_active_sector_total_dim}
    \end{align}
    Thus the dimension is the total dimension of the active blocks, not the
    ambient bond dimension $D$. This statement records positive-length MPV
    equality; the exact coisometric reconstruction is stated next.
\end{theorem}

\begin{proof}\leanok
    \uses{thm:cpsv_active_sector_bnt_exact, thm:coisometry_reconstruction_same_mpv_pos, thm:gauge_invariance}
    Apply Theorem~\ref{thm:cpsv_active_sector_bnt_exact}. Its reconstruction
    writes $A^i=U^\dagger B^iU$, where $UU^\dagger=\Id$ and
    $B^i=XP^iX^{-1}$. Theorem~\ref{thm:coisometry_reconstruction_same_mpv_pos}
    shows that $A$ and $B$ generate the same matrix product vectors at every
    positive length. Theorem~\ref{thm:gauge_invariance} gives the same conclusion
    for $B$ and $P$. Transitivity yields $\mc{V}^+(A)=\mc{V}^+(P)$, while the
    dimension identity is already part of the exact reconstruction theorem.
\end{proof}

% Source: CPSV16 lines 214--301 and 1080--1148; CPGSV21 lines 1831--1885.
\begin{theorem}[Exact active coisometric BNT reconstruction]%
    \label{thm:cpsv_active_sector_bnt_exact}
    \lean{MPSTensor.CPSVCanonicalFormData.exists_active_isBNTCanonicalForm_exact}
    \leanok
    \uses{def:cpsv_canonical_form, def:sector_bnt_cf}
    Let $A$ have CPSV canonical-form data satisfying the separate normalization
    condition~(\ref{eq:cpsv_weight_normalization}), assume $A\ne0$, and put
    $K_{\mathrm{act}}=\{k:\mu_k\ne0\}$. There are a BNT canonical-form sector
    decomposition $P$, a coisometry
    $U\in\C^{D_{\mathrm{tot}}(P)\times D}$, and
    $X\in\GL_{D_{\mathrm{tot}}(P)}(\C)$. The first two identities below hold,
    and the third holds for every physical index $i$:
    \begin{align}
        D_{\mathrm{tot}}(P)
                   & =\sum_{k\in K_{\mathrm{act}}}D_k,
        \label{eq:cpsv_active_exact_bnt_dim}             \\
        UU^\dagger & =\Id,
        \label{eq:cpsv_active_exact_bnt_coisometry}      \\
        A^i        & =U^\dagger\left(XP^iX^{-1}\right)U.
        \label{eq:cpsv_active_exact_bnt_reconstruction}
    \end{align}
\end{theorem}

\begin{proof}\leanok
    \uses{thm:cpsv_exact_active_reconstruction, thm:cpsv_normal_tp_gauge, lem:gauge_equiv_from_block_gauge, thm:paperbnt_supplier_prepared_exact}
    By Theorem~\ref{thm:cpsv_exact_active_reconstruction}, the original tensor
    has the exact active reconstruction
    \begin{align}
        A^i        & =V^\dagger T_{\mathrm{act}}^iV,
        \label{eq:cpsv_active_exact_bnt_start}       \\
        VV^\dagger & =\Id.
        \label{eq:cpsv_active_exact_bnt_start_coisometry}
    \end{align}
    Put every active normal block into a trace-preserving primitive irreducible
    gauge $Y$, and apply
    Theorem~\ref{thm:paperbnt_supplier_prepared_exact}. Let $Q$ be the
    permutation matrix of its flattened coordinate equivalence, so
    $QQ^\dagger=\Id$. Reindexing $Y$ by the same equivalence gives an
    invertible matrix $W$ on the bond space of $P$ and the identities
    \begin{align}
        QY              & =WQ,
        \label{eq:cpsv_active_exact_bnt_intertwining} \\
        Y^{-1}Q^\dagger & =Q^\dagger W^{-1}.
        \label{eq:cpsv_active_exact_bnt_intertwining_inv}
    \end{align}
    If $X_0$ is the prepared-block gauge, set
    $X=W^{-1}X_0$ and $U=QV$. Substitution in
    (\ref{eq:cpsv_active_exact_bnt_start}) gives
    (\ref{eq:cpsv_active_exact_bnt_reconstruction}), while
    \begin{align}
        UU^\dagger
         & =Q(VV^\dagger)Q^\dagger=\Id.
        \notag
    \end{align}
    This proves (\ref{eq:cpsv_active_exact_bnt_coisometry}). The active
    enumeration gives (\ref{eq:cpsv_active_exact_bnt_dim}).
\end{proof}

% Source anchors: RMP arXiv:2011.12127v2, lines 1831--1906;
% CPSV16 arXiv:1606.00608, lines 354--361 and 1172--1199.
\begin{theorem}[Equal fundamental theorem for active CPSV canonical forms]%
    \label{thm:cpsv_active_core_equal_ft}
    \lean{MPSTensor.CPSVCanonicalFormData.exists_active_fundamentalTheorem_equal_canonicalForm}
    \leanok
    \uses{thm:cpsv_active_core_equal_ft_unitary}
    Let $A$ and $B$ have CPSV canonical-form data with weights
    $(\mu_k^A)_k$, $(\mu_k^B)_k$ and block dimensions
    $(D_k^A)_k$, $(D_k^B)_k$. Assume that both data satisfy the separate
    normalization condition~(\ref{eq:cpsv_weight_normalization}), that
    $A\ne0$ and $B\ne0$, and that
    $\mc{V}^+(A)=\mc{V}^+(B)$. Put
    $K_A=\{k:\mu_k^A\ne0\}$ and $K_B=\{k:\mu_k^B\ne0\}$.
    Then there are BNT canonical-form sector decompositions $P$ and $Q$ such
    that
    \begin{align}
        \mc{V}^+(A)          & =\mc{V}^+(P),
        \notag                                        \\
        \mc{V}^+(B)          & =\mc{V}^+(Q),
        \notag                                        \\
        D_{\mathrm{tot}}(P)  & =\sum_{k\in K_A}D_k^A,
        \notag                                        \\
        D_{\mathrm{tot}}(Q)  & =\sum_{k\in K_B}D_k^B,
        \notag                                        \\
        D_{\mathrm{tot}}(P)  & =D_{\mathrm{tot}}(Q),
        \notag                                        \\
        \sum_{k\in K_A}D_k^A & =\sum_{k\in K_B}D_k^B.
        \label{eq:cpsv_active_core_equal_dims}
    \end{align}
    After identifying the two total bond spaces by their equality, there is
    an invertible matrix $X$ such that, for every physical index $i$,
    \begin{align}
        Q^i & =X P^iX^{-1}.
        \label{eq:cpsv_active_core_equal_gauge}
    \end{align}
    The conclusion concerns the extracted active tensors $P$ and $Q$. It does
    not assert equality of the ambient bond dimensions $D_A$ and $D_B$, and it
    does not assert that $A$ and $B$ are gauge equivalent.
    % Scope restriction (active canonical-form tensors): active-only corrected form;
    % see docs/paper-gaps/tnlean_bnt_ft_theorem_surface.tex.
\end{theorem}

\begin{proof}\leanok
    \uses{thm:cpsv_active_core_equal_ft_unitary}
    Apply Theorem~\ref{thm:cpsv_active_core_equal_ft_unitary} and retain its
    extracted sector decompositions, dimension identities, and conjugating
    matrix. Forget that the conjugating matrix is unitary.
\end{proof}

% The RMP unitary note is arXiv:2011.12127v2, lines 1905--1906.
% The CPSV16 CFII refinement is arXiv:1606.00608, lines 1197--1199.
\begin{theorem}[Unitary refinement for active CPSV canonical forms]%
    \label{thm:cpsv_active_core_equal_ft_unitary}
    \lean{MPSTensor.CPSVCanonicalFormData.exists_active_fundamentalTheorem_equal_canonicalForm_unitary}
    \leanok
    \uses{thm:cpsv_active_sector_bnt, cor:sector_bnt_equal_unitary_global_gauge}
    Let $A$ and $B$ have CPSV canonical-form data with weights
    $(\mu_k^A)_k$, $(\mu_k^B)_k$ and block dimensions
    $(D_k^A)_k$, $(D_k^B)_k$. Assume that both data satisfy the separate
    normalization condition~(\ref{eq:cpsv_weight_normalization}), that
    $A\ne0$ and $B\ne0$, and that
    $\mc{V}^+(A)=\mc{V}^+(B)$. Put
    $K_A=\{k:\mu_k^A\ne0\}$ and $K_B=\{k:\mu_k^B\ne0\}$.
    Then there are BNT canonical-form sector decompositions $P$ and $Q$ such
    that
    \begin{align}
        \mc{V}^+(A)          & =\mc{V}^+(P),
        \notag                                        \\
        \mc{V}^+(B)          & =\mc{V}^+(Q),
        \notag                                        \\
        D_{\mathrm{tot}}(P)  & =\sum_{k\in K_A}D_k^A,
        \notag                                        \\
        D_{\mathrm{tot}}(Q)  & =\sum_{k\in K_B}D_k^B,
        \notag                                        \\
        D_{\mathrm{tot}}(P)  & =D_{\mathrm{tot}}(Q),
        \notag                                        \\
        \sum_{k\in K_A}D_k^A & =\sum_{k\in K_B}D_k^B.
        \label{eq:cpsv_active_core_unitary_dims}
    \end{align}
    After identifying the two total bond spaces, there is a unitary matrix $U$
    such that, for every physical index $i$,
    \begin{align}
        Q^i & =U P^iU^\dagger.
        \label{eq:cpsv_active_core_unitary_gauge}
    \end{align}
    This is the unitary choice noted in
    \cite[lines~1905--1906]{Cirac2021Matrix} and the canonical-form-II
    refinement of \cite[Corollary~A.6, lines~1197--1199]{Cirac2016MPDO_arXiv}.
    Again, there is no claim that $D_A=D_B$, and no gauge equivalence between
    $A$ and $B$ is asserted.
    % Scope restriction (active canonical-form tensors): active-only corrected form;
    % see docs/paper-gaps/tnlean_bnt_ft_theorem_surface.tex.
\end{theorem}

\begin{proof}\leanok
    \uses{thm:cpsv_active_sector_bnt, cor:sector_bnt_equal_unitary_global_gauge}
    Apply Theorem~\ref{thm:cpsv_active_sector_bnt} to $A$ and $B$. This gives
    $P$ and $Q$, their BNT canonical-form properties, the first four identities
    in~(\ref{eq:cpsv_active_core_unitary_dims}), and
    \begin{align}
        \mc{V}^+(P) & =\mc{V}^+(Q).
        \notag
    \end{align}
    Corollary~\ref{cor:sector_bnt_equal_unitary_global_gauge} gives the fifth
    identity and supplies $U$. Substituting the two active
    dimension formulas gives the last equality in
    (\ref{eq:cpsv_active_core_unitary_dims}).
\end{proof}

% Source anchors: RMP arXiv:2011.12127v2, lines 1831--1906;
% CPSV16 arXiv:1606.00608, lines 237--301 and 1135--1199.
\begin{theorem}[Equal-ambient gauge equivalence for CPSV canonical forms]%
    \label{thm:cpsv_equal_ambient_gauge}
    \lean{MPSTensor.CPSVCanonicalFormData.%
        fundamentalTheorem_equal_ambient_canonicalForm}
    \leanok
    \uses{def:cpsv_canonical_form, def:same_mpv2_pos, def:gauge_equiv}
    Let $A$ and $B$ have the same ambient bond dimension $D$. Suppose each
    tensor has CPSV canonical-form data, both data satisfy the separate weight
    normalization condition~(\ref{eq:cpsv_weight_normalization}), neither tensor
    is zero, and $\mc{V}^+(A)=\mc{V}^+(B)$. Then $A$ and $B$ are gauge equivalent:
    there is
    $G\in\GL_D(\C)$ such that, for every physical index $i$,
    \begin{align}
        B^i & =GA^iG^{-1}.
        \label{eq:cpsv_equal_ambient_gauge}
    \end{align}
    Zero-weight listed blocks and unused ambient coordinates are allowed. The
    conclusion requires the common ambient dimension $D$ and makes no claim for
    canonical forms with different ambient bond dimensions.
\end{theorem}

The cited sources supply the active canonical-form comparison. The passage back
from the active tensors to the ambient tensors is the coisometric lifting of
Theorem~\ref{thm:coisometry_reconstruction_gauge_equiv}.

\begin{proof}\leanok
    \uses{thm:cpsv_active_sector_bnt_exact, thm:coisometry_reconstruction_same_mpv_pos, thm:cpgsv_equal_case_source, thm:coisometry_reconstruction_gauge_equiv, thm:gauge_invariance}
    Apply Theorem~\ref{thm:cpsv_active_sector_bnt_exact} to obtain
    BNT canonical forms $P,Q$, coisometries $U_A,U_B$, and retained-space
    gauges $X_A,X_B$ with
    \begin{align}
        A^i & =U_A^\dagger X_AP^iX_A^{-1}U_A,
        \notag                                \\
        B^i & =U_B^\dagger X_BQ^iX_B^{-1}U_B.
        \label{eq:cpsv_equal_ambient_reconstructions}
    \end{align}
    Theorem~\ref{thm:coisometry_reconstruction_same_mpv_pos}, gauge invariance,
    and the hypothesis give
    \begin{align}
        \mc{V}^+(P)
         & =\mc{V}^+(A)=\mc{V}^+(B)=\mc{V}^+(Q).
        \label{eq:cpsv_equal_ambient_active_mpv}
    \end{align}
    Theorem~\ref{thm:cpgsv_equal_case_source} therefore identifies the retained
    dimensions and gives $Y$ such that $Q^i=YP^iY^{-1}$. After identifying the
    retained spaces, the two middle tensors in
    (\ref{eq:cpsv_equal_ambient_reconstructions}) are conjugate by
    $X_BYX_A^{-1}$. Theorem~\ref{thm:coisometry_reconstruction_gauge_equiv}
    aligns the two coisometries inside the common ambient space and yields
    (\ref{eq:cpsv_equal_ambient_gauge}).
\end{proof}

Combining the after-blocking preparation with the normalized-weight construction
gives the arbitrary-input statement, conditional on the weight normalization.

\begin{remark}[Arbitrary-input reduction]%
    \label{rem:arbitrary_input_reduction_gap}
    The reduction is conditional on the two weight-normalization hypotheses of
    Theorem~\ref{thm:paperbnt_supplier_after_blocking}.
    % Local fix (cpsv16_cf_normalization_and_proportional_comparison): the
    % preparatory blocks (thm:prepared_bnt_blocks_after_blocking) carry no weight
    % normalization. The remaining step rescales
    % the weights to |\mu_k| <= 1 with a unit witness, recorded in
    % docs/paper-gaps/cpsv16_cf_normalization_and_proportional_comparison.tex.
    The weight rescaling is carried out in
    Theorem~\ref{thm:paperbnt_supplier_after_blocking_normalized} below, which
    discharges both hypotheses at the price of one positive scalar per site.
\end{remark}

\begin{theorem}[Normalized BNT form after blocking for a nonzero MPV family]%
    \label{thm:paperbnt_supplier_after_blocking_normalized}
    \lean{MPSTensor.exists_isBNTCanonicalForm_afterBlocking_pos_normalized}
    \leanok
    \uses{thm:prepared_bnt_blocks_after_blocking, thm:paperbnt_supplier_prepared, thm:paperbnt_weight_normalization, lem:mpv_to_tensor_from_blocks_weight_mul_left, thm:mpv_decomposition, lem:eval_word_append, lem:eval_word_block}
    Let $A$ be an MPS tensor whose matrix-product family has a nonzero
    coefficient at some positive length. Then there exist a blocking length
    $p\geq1$, a positive real $m$, and a sector decomposition $P$ over the
    blocked alphabet in the BNT canonical form of
    Definition~\ref{def:sector_bnt_cf}, such that for every positive length
    $N$,
    \begin{align}
        V^{(N)}(A^{[p]})
         & =m^{N}\,V^{(N)}(P).
        \label{eq:bnt_sector_normalized_family}
    \end{align}
    The scalar $m^{N}$ is the per-site normalization cost of the line-246
    weight choice; the hypotheses of
    Theorem~\ref{thm:paperbnt_supplier_after_blocking} are discharged rather
    than assumed.

    \textbf{Scope restriction (nonzero MPV family).} The source proposition
    says ``for any tensor'', but the zero MPV family cannot supply the global
    unit-modulus weight imposed immediately before it. The hypothesis above
    is exactly the exclusion of this degenerate case. The normalization issue
    and the proportional-MPV comparison are documented in
    \path{docs/paper-gaps/cpsv16_cf_normalization_and_proportional_comparison.tex}.
\end{theorem}

\begin{proof}\leanok
    Choose a word $w$ with $\tr(A^w)\neq0$. For every $p\geq1$, there is a
    $q\geq1$ such that $\tr((A^w)^{pq})\neq0$. Otherwise Newton--Girard applied
    to $(A^w)^p$ and the zero matrix would show that $(A^w)^p$, and hence
    $A^w$, is nilpotent, contradicting $\tr(A^w)\neq0$. Repeating $w$ exactly
    $pq$ times gives a nonzero coefficient of the $p$-blocked tensor. Thus the
    prepared family at the blocking length supplied by
    Theorem~\ref{thm:prepared_bnt_blocks_after_blocking} is nonempty.

    Normalize its weights by their maximal modulus $m$
    (Theorem~\ref{thm:paperbnt_weight_normalization}). By the weight-scaling
    identity of
    Lemma~\ref{lem:mpv_to_tensor_from_blocks_weight_mul_left},
    \begin{align}
        V^{(N)}(A^{[p]})
         & =
        V^{(N)}\left(\bigoplus_k \mu_k\,B_k\right)
        =
        V^{(N)}\Bigl(\bigoplus_k
        m\,\frac{\mu_k}{m}\,B_k\Bigr),
        \notag \\
         & =
        m^{N}\,V^{(N)}\left(\bigoplus_k \frac{\mu_k}{m}\,B_k\right)
        =
        m^{N}\,V^{(N)}(P).
        \notag
    \end{align}
    where the last identity is the prepared supplier of
    Theorem~\ref{thm:paperbnt_supplier_prepared} applied to the normalized
    weights, which gives the BNT canonical form $P$.
\end{proof}
